% Source: ACM_Templates.pdf pp. 40-66. pp. 55-57; https://cp-algorithms.com/data_structures/sqrt_decomposition.html
\subsection{普通莫队：区间不同数个数}
\label{ds:mo}
\AlgorithmInfo
    {离线处理静态区间统计，示例维护区间不同值数量。}
    {块长 $B$：移动 $O(n^2/B+qB+n)$，排序 $O(q\log(q+1))$；空间 $O(n+m+q)$。}
    {数组 1 下标，值先离散化为 $1\ldots m$；查询为非空闭区间，不含修改。}
\subsubsection*{原理}
若区间左右端点各移动一步时能快速增删元素，可以离线重排询问以减少总移动量。先按左端点所在块排序，同块内右端点交替升降。维护当前窗口 $[l,r]$、各值出现次数与不同值数，扩张窗口时加入元素，缩小时删除。

块长取约 $n/\sqrt q$，至少为 1；$n,q$ 同阶时移动量为 $O(n\sqrt n)$。先扩张再收缩，避免在两个不相交窗口之间删除尚未加入的元素。

\subsubsection*{使用教程}
\InlineCode{init(a,m)} 保存原数组并清空询问；\InlineCode{add_query(l,r)} 追加询问并返回从 0 开始的编号。\InlineCode{work()} 返回按加入顺序排列的答案，不覆盖原数组与询问，可重复调用；没有询问时返回空数组。更换统计内容时修改局部 \InlineCode{add/del} 和当前答案状态。
\begin{lstlisting}[style=mtycpp,numbers=none]
Mo mo;
mo.init(vector<int>{0,1,2,1,3,2},3);
int id=mo.add_query(1,3);
mo.add_query(2,5);
auto ans=mo.work();
int kinds=ans[id];
\end{lstlisting}
\needspace{7\baselineskip}
\subsubsection*{模板代码}
\CodeFile{code/data-structures/mo.cpp}
\subsubsection*{需要注意}
值域包含负数或很大时先离散化；频次数组按值编号访问，不按位置访问。本例添加一个元素和删除一个元素均为 $O(1)$，替换统计量后应重新计入维护代价。
